\needspace{5cm} \section{Divák} \label{divák} Divák je nedílnou součástí veřejných představení. 
Nejen že platí vstupné, a pak svým smíchem a potleskem poskytuje představení palivo/energii; svými příběhy poskytuje úvodní impuls k~delším formám, může si sám vyzkoušet mírné zapojení od \odkaz{žánrové vyprávěné}{žánrová vyprávěná} po \odkaz{Love song}{love song} či \odkaz{Moving people}{moving people}.   V~některých \odkaz{formách}{forma} také uděluje body. 
Následně ideálně divák přijde příště zas, a ideálně i přivede další kamarády. 
 
 
\needspace{5cm} \section{Hráč} \label{hráč} Označení pro jedince na scéně, kteří jsou zodpovědní za tvorbu obsahu improvizace. Mezi hráče nepatří \odkaz{MC}{mc} a \odkaz{hudebníci}{hudebník}. V~rámci improvizačních představení se preferuje výraz "\textbf{hráč}\uv{ nad výrazem } \textbf{herec}"{} i když mají téměř stejný význam. Všichni jedinci na scéně jsou herci, ale hráči jsou jedinci zastupující vždy svůj tým (zápas) či jsou vyčleněni od \odkaz{MC}{mc} (improshow).  
 
\subsection{ Zápasová forma } Na \odkaz{zápasech}{zápas} jsou hráči rozděleni do \odkaz{týmů}{tým}. 
 
\subsection{ Zápasový oděv } Na veřejnou rozcvičku jsou hráči oděni čistě v~černém (bez výšivek, potisků apod),  
na samotný zápas si pak berou trika v~barvách svých týmů. 
Trikot tedy zahrnuje 
\begin{itemize}
\item ponožky a kalhoty   případně leginy/punčocháče (40 DEN a víc) a sukni v~barvě černé
\item triko v~barvě černé
\item triko v~barvě týmu.
\end{itemize}
V~případě dlouhých vlasů je takřka nutnost gumička do vlasů, jinak nebývá skrz bujnou hřívu vidět do obličeje. 
 
 
 
 
 
\needspace{5cm} \section{Hudebník} \label{hudebník} Hudebník je skvělým oživením každého zápasu či improshow. 
Jeho role je proměnlivá, při \odkaz{hlasování}{hlasování} na zápasech hraje výrazně až freneticky. 
Ve všech \odkaz{zpívaných kategoriích}{:kategorie:zpívané kategorie} je určovatelem/hybatelem charakteru písní. 
A~v~dalších kategoriích může podkreslovat. 
 
Skvělé je mít několik nástrojů (či klávesy s~přepínáním rejstříků) a volit dle charakteru děje/ hlasu hráčů, či jen tak pro zpestření. 
Je vhodné, když, zejména pro méně hudebně zkušené hráče, zvýrazní první dobu a případně i ukáže kývnutím. 
 
\subsection{Role hudebníka} Na hudebníka může být mnoho různých požadavků. Záleží na oraganizátorovi představení, formátu představení, zkušenostech hudebníka,  zkušenostech herců a pod. Jeho role může být: 
\begin{itemize}
\item \textbf{Hudební podkres mimo improvizaci} - Před začátkem představení, během úvodu, vyhlášení pauzy, hlasování, ukončení přestavení. Žánr záleží na cílené atmosféře celého představení, ale obvykle by měla být hudba pozitivní, příjemná, veselá až radostná. Pokud jde o~jasné oznámení začátku nebo konce, je vhodná i fanfára.
\item \textbf{Podkresová hudba} - Slouží pro doplnění nálady a atmosféry během improvizace. Neměla by brát \odkaz{fokus}{fokus} hercům.
\item \textbf{Krátké znělky (jingly)} - Mezi větami či souslovími, například v~kategoriích \odkaz{1000 způsobů jak}{1000 způsobů jak}, \odkaz{Věty, které by neměly zaznít}{věty, které by neměly zaznít} nebo v~rozcvičce \odkaz{Jsem, beru}{jsem, beru}.
\item \textbf{Doprovod písně} - Píseň by měla začít jasnou předehrou, ve které hudebník deklaruje tóninu, (jednoduchý) rytmus, tempo a ideálně i (jednoduchý) akordový postup, kterého se po zbytek písně více méně drží. Možnosti variací jsou velmi závislé na schopnostech herce a na jeho souhře s~hudebníkem. Experimentování se nedoporučuje, pokud si hudebník není opravdu jistý hercovými schopnostmi.
\item \textbf{Zvukové efekty} - Klepání na dveře, vyzvánení telefonu nebo třeba pípající přístroj v~nemocnici improvizaci okoření.
\item \textbf{Ukončení improvizace} - Improvizaci může ukončit i hudebník. Zvláště, pokud byla právě řečena skvělá pointa na závěr, je to chvíle, kdy dát hercům případně \odkaz{MC}{mc} hudbou najevo, že je vhodný okamžik improvizaci ukončit. V~písních má hudebník na starosti ukončování ve většině případů.
\item \textbf{Nový impulz do improvizace} - Ve vhodné chvíli, může hudebník výrazně změnit styl doprovodu a přinést tak do improvizace impulz, kterého se herci chytí. Je důležité tuto změnu udělat tak výrazně, aby si toho herci všimli. Vhodná situace je například ve chvíli, kdy se v~příběhu dlouho nic neděje.
\item \textbf{Umět nehrát} - Do některých kategorií se hudba hodí více, do některých méně. Vždy je ale důležité nezapomínat na využití ticha a na přehlcení hudbou.
\end{itemize}
 
\subsection{Schopnosti a dovednosti} \begin{itemize}
\item \textbf{Poznat, kdy hrát a kdy nehrát}
\item \textbf{Hrát v~\odkaz{hudební nálady}{náladě}, která se hodí} nebo aspoň neodporuje
\item \textbf{Umět improvizovat.} Naučený repertoir má své limity.
\item \textbf{Umět zahrát v~mnoha náladách}, nejen \textit{veselý - smutný}
\item \textbf{Umět zahrát v~mnoha hudebních žánrech} a doprovodit mnoho filmových a literárních žánrů
\item \textbf{Umět zahrát improvizovanou píseň}, která je čitelná pro herce
\item \textbf{Umět pomoci hercům dodržet rytmus, tóninu}, vrátit se zpátky do rytmu, tóniny, nebo se podle herců přizpůsobit.
\item \textbf{Neplést herce složitostmi}
\end{itemize}
 
\subsection{Kdy hrát a kdy nehrát} Některé kategorie hudbu přímo vyžadují, v~některých by hudba překážela. Je dobré si hudební doprovod dvakrát rozmyslet v~kategoriích, které jsou založené na zvucích odjinud, jako \odkaz{Zvuky}{zvuky} nebo \odkaz{Dabing filmu}{dabing filmu}. Dále se hudba příliš nehodí do kategorií s~rychlým spádem (\odkaz{Schizofrenie}{schizofrenie}) a obecně kategorií náročných na pozornost (\odkaz{Sportovní komentátor}{sportovní komentátor}). Do kategorií založených na gagování (\odkaz{1000 způsobů jak}{1000 způsobů jak}, \odkaz{Metafory}{metafory}) se hodí pouze oddělování znělkou (jinglem), do zpívaných kategorií je zase vhodné nehrát mezi písničkami, aby nebylo hudby příliš. Hudba se naopak hodí pro kategorie, které vybízejí ke dlouhému příběhu, ke kategoriím, kde se střídají prostředí (\odkaz{Poslední věta}{poslední věta}) nebo úhly pohledů (\odkaz{Spoon river}{spoon river}) pro jasné odlišení. Hudební podkres se dá použít ke gradaci v~kategoriích, které jinak hudbu příliš nevyžadují (\odkaz{Polovina času}{polovina času}, \odkaz{Smrt v~1 minutě}{smrt v~1 minutě}). Hudbou se dá vyplnit prázdný prostor nebo naopak podtrhnout silný okamžik, je ale dobré rozvrhnout si množství hudby do celého představení, aby hudba nehrála pořád. V~\odkaz{Longformách}{longforma} má hudba na starosti úvod, závěr, předěly a má podkreslovat důležité okamžiky, ale zároveň má nechat prostor na většinu dialogů pouze hercům.  
 
\subsection{Spolupráce s~hráči} Hráč může před představením probrat s~hudebníkem své představy - např. preference písní pomalých, od C apod. 
 
Zejména běžné bývá upozornění hudebníka, že míváte problém chytit rytmus, hudebník pak na vás (i dle hudebního nástroje) kývne, či zařadí před začátek fráze, ve které by bylo záhodno začít zpívat, drobnou pauzu. Pravděpodobně také zvolí ve vašem případě nějaký jednodušší hudební styl na 4 doby. 
 
 
\needspace{5cm} \section{Konferenciér} \label{konferenciér} Konferenciér řídí průběh \odkaz{zápasu}{zápas} či \odkaz{improshow}{improshow}. 
 
\subsection{Běžné chyby} \begin{itemize}
\item rozvleklý úvod
\item chybný či nepřehledný popis kategorií
\end{itemize}
 
\subsection{ Kam dál } \begin{itemize}
\item  \odkaz{Manuál konferenciéra}{manuál konferenciéra}
\end{itemize}
 
 
 
 
\needspace{5cm} \section{MC} \label{mc} \textbf{MC} z~anglického \textbf{Master of Ceremony}. \odkaz{Konferenciér}{konferenciér}/moderátor/\odkaz{rozhodčí}{rozhodčí}. Prostě ten, kdo má pod pantoflem celý průběh \odkaz{zápasu}{zápas}/\odkaz{improshow}{improshow}. 
Na improshow někdy bývá \uv{kolující MC}  - \odkaz{hráči}{hráč} se v~řízení průběhu střídají. 
 
 
\needspace{5cm} \section{Osvětlovač} \label{osvětlovač} Osvětlovač se stará, aby bylo vidět to, co má být vidět; a naopak. 
 
 
U~\odkaz{zápasu}{zápas} rozsvěcí a zhasíná, či tlumí světla zejména na pokyny od \odkaz{rozhodčího}{rozhodčí} (např. Můžeme hlasovat. ).  
U~\odkaz{improshow}{improshow} či \odkaz{longforem}{longforma} může být častěji hybatelem děje, zejména formou střihů/ukončování scén. 
 
Rozhodně se může zapojit volbou barvy a intenzity světla, nebo dynamickými změnami  (Disko, bouřka). 
 
 
\subsection{Osvětlovací pult} Typicky sedí u~osvětlovacího pultu, který digitálně ( DMX 512 sběrnice procházející skrze všechny světelné zdroje/mlhovače apod) nebo analogově (výkonové stmívací zásuvky)  ovládá všechna nebo většinu světel v~sále. 
Vyjímkou může tvořit pracovní světlo na scéně (pracák) a v~sále ovládané klasickými spínači na místě na stěně. 
 
Osvětlovací pult obsahuje řadu \uv{šavlí}  - tahových potenciometrů, kterými může plynule měnit intenzitu světla jednotlivých světel, skupin světel či barevných kanálů. 
Zároveň bývá na pultu jedna master šavle, kterou je možno snížit intenzitu všeho.  
Pod  šavlemi nalezneme spínač, kterým můžeme kanál okamžitě vypnout (např. již zmiňovaná bouřka). 
 
Je vhodné si ještě před představením pult a rozmístění světel vyzkoušet, rozmyslet si, která světla se na co hodí, a případně je přestavět/doplnit filtry. 
 
 
 
 
 
\subsection{ Barvy světla } \subsubsection{ 	Bílá } Základní a neutrální, při zvýšené intenzitě působí sterilně, uměle. 
Všechny typy scén, velmi vhodné do nemocničního, vědeckého a kancelářského prostředí. 
\subsubsection{ 	Žlutá } Může nahrazovat bílé světlo, skvěle se hodí na venkovní denní osvětlení. Při zvýšené intenzitě vytváří efekt sucha a horka. 
Všechny typy scén, zvláště pak ty venkovní. Vhodné na dokreslení letního období. Intenzivní žluté světlo se hodí pro pouště, western a cizí planety. 
\subsubsection{ 	Červená} Nepřirozené, silně příznakové světlo. Pojí se s~vášní, láskou, zuřivostí a násilím. 
Vhodné pro detektivky/kriminálky (popř. pro noirové scény), horory a pro červenou knihovnu. V~technickém a sci-fi prostředí skvěle dokresluje různé havárie a poruchy. Při nižší intenzitě může vytvářet efekt večerních červánků. 
Náhlé zrudnutí scény může znamenat krev, krvavou smrt, sex. 
\subsubsection{ 	Modrá } „Světlo pro chlad, smutek, vodu a tmu“.  
Vhodné pro všechny druhy nočních a večerních scén, pro zimu a polární prostředí. Hodí se k~tragickým scénám (bez krve). Při nízké intenzitě (v~kombinaci se silnou bílou) velmi dobře dokresluje sterilní prostředí. Ideální je pro podvodní scény a pro deštivé počasí.  
\subsubsection{ 	Zelená } Silně nepřirozené světlo (s~výjimkou scén z~lesního prostředí), málokdy „ladí“ s~barvou kůže a vlasů hráčů/herců (tluče se především s~odstíny červené a růžové, hnědou, oranžovou a žlutou přebíjí). Vyžaduje velmi citlivé zacházení osvětlovače. 
Skvěle funguje ve scénách, které se odehrávají v~lese, v~džungli, v~bažině. 
Díky své nepřirozenosti vhodně dokresluje zapáchající skládky, nadpřirozené jevy/bytosti a radiaci (především ve sci-fi a postapokalyptických scénách). 
 
 
\subsection{	Míchání barev } Přidáváním barevného světla do bílého (popř. do žlutého) vzniká mnohem jemnější a přirozenější osvětlení. Intenzivní čistě barevné světlo je silně příznakové.  
\subsubsection{ Příklady míchání barev s~možným využitím } \begin{itemize}
\item  Bílá + žlutá → nejpřirozenější světlo
\item  Bílá + červená → romantika, červánky, něžnosti
\item  Bílá + modrá → zima, chlad, sterilní prostředí (nemocnice apod.)
\item  Žlutá + červená → horko, léto, požáry
\item  Žlutá + (méně intenzivní) modrá → les, přímořské oblasti, pláže
\item  Žlutá + zelená → přirozené lesní světlo (lepší než čistě zelené), postapokalyptické výjevy
\item  Červená + modrá → magie, večer, červánky, polární záře, tajemno
\item  Červená + zelená → technická havárie, únik radiace, pralesní nebezpečí (např. lidožrouti, masožravé rostliny apod.)
\item  Modrá + zelená → voda  (více modré → čerstvá, čistá nebo slaná voda; více zelené → stojatá, špinavá a rybniční voda), polární záře, nadpřirozeno
\end{itemize}
 
 
	Pro dosažení efektu pulzujícího či proměnlivého světla (oheň, polární záře, světlo pronikající skrz koruny stromů apod.) je vhodné průběžně zvyšovat a snižovat intenzitu jednotlivých barevných světel (např. u~požáru nezávisle na sobě měnit intenzitu žlutého a červeného světla).  
	Osvětlovač by se neměl bát si se světly hrát a objevovat nové využití jednotlivých barev. 
 
\subsection{ Intenzita } \begin{itemize}
\item  Intenzita světla by měla odpovídat typu scény – například ve strašidelném hradě očekáváme spíš šero a tlumené světlo, na operačním sále je na místě ostré bílé světlo apod.
\item  Intenzitou světla lze snadno dokreslit denní dobu. Jinak ozářená bude scéna v~noci, za svítání a v~poledne.
\item  Různou intenzitou bodových světel se určuje důležité místo (popř. hráč/herec) na jevišti.
\end{itemize}
 
 
\subsection{ Směr a rozptyl světla } \begin{itemize}
\item 	\textbf{Rozptýlené}  Neutrální, plní základní funkci osvětlení jeviště.
\item 	\textbf{Bodové shora} 	Spotlight. Jasně určuje, co je v~dané chvíli na jevišti nejdůležitější. Vhodné pro monology a sólové pěvecké či recitační výstupy.
\item 	\textbf{Ze strany} Vizuálně zajímavé osvětlení – vytváří na objektech a na hercích/hráčích mnoho stínů a kontrastů. Velmi vhodné pro napínavé nebo vyostřené scény.
\item 	\textbf{Zepředu na oponu} 	Při zatažené oponě buduje napětí (divák tuší, že se co nevidět něco začne dít). Může sloužit jako spotlight pro konferciéra, pro rozhodčího i pro hráče/herce.
\item 	\textbf{Nad diváky} 	Je rozsvícené před a po představení a během přestávky. Během představení je rozsvícené (alespoň při nižší intenzitě) v~případech, kdy někdo z~účinkujících (zpravidla konferenciér či zápasový rozhodčí) přímo interaguje s~diváky.
\end{itemize}
 
Na co pozor? 
Osvětlovač by si měl dát pozor na to, aby byli vždy osvětleni všichni hráči/herci, kteří mají být vidět (např. aby při dialogu neosvětloval jen jednoho z~účinkujících kvůli „hezkému efektu“). 
 
\subsection{Další využití barev}  
V~případě, že se v~představení opakuje omezený počet pevně daných prostředí (např. v~kategorii čtverec), může osvětlovač každému z~těchto prostředí přidělit jednu barvu či kombinaci barev. Je to způsob, jakým může divákům (a možná i hráčům/hercům) usnadnit orientaci v~různých prostředích. Vyžaduje to však velké soustředění a přesnost ze strany osvětlovače (aby rozsvícením špatného světla nevyvolal naopak zmatek). 
 
Podobně může osvětlovač přiřadit konkrétní barvu (ovšem spíš v~nižší intenzitě)  jednotlivým postavám (takový přístup je vhodný spíše u~longforem, ve kterých vystupují po celou dobu představení stejné postavy). Pokud jsou na scéně současně dvě postavy, může jejich barvy zkombinovat. V~případě, že je na jevišti více postav, může se osvětlovač uchýlit k~neutrálnímu bílému světlu, popř. k~barvě náležící nejvýraznější postavě. Také tento přístup si vyžaduje osvětlovačovu plnou soustředěnost. 
 
 
 
\needspace{5cm} \section{Pomocný rozhodčí} \label{pomocný rozhodčí} \textbf{Pomocný rozhodčí}, též \textbf{pomocňák} je typicky k~vidění na \odkaz{zápasech v~improvizaci}{zápas}. 
Bývají dva, jejich úkolem je \odkaz{měřit čas}{ukazování času} kategorií, sčítání \odkaz{hlasů}{hlasování} (\odkaz{kartičky}{hlasovací kartička}), aktualizace skóre a faulů, podávání košíčku s~tématy. Jejich pomoc \odkaz{hlavnímu rozhodčímu}{rozhodčí} někdy může zahrnovat zapisování a sledování \odkaz{faulů}{faul}, které rozhodčí pískne.  
 
Pomoci vám může i \odkaz{Manuál pomocného rozhodčího}{manuál pomocného rozhodčího}. 
\subsection{Běžné chyby} \begin{itemize}
\item Pomocňáci si nedohodnou, který sčítá hlasy kterému týmu a sčítání (typicky první na zápase, či první po přestávce) se musí opakovat.
\end{itemize}
 
 
 
\needspace{5cm} \section{Rozhodčí} \label{rozhodčí} \textbf{Rozhodčí} je klíčovou postavou improvizačních \odkaz{zápasů}{zápas}, obvykle tou s~nejvyším \odkaz{statusem}{status} na scéně. 
Během zápasu zejména vybírá \odkaz{kategorie}{kategorie}, píská \odkaz{fauly}{faul}, uděluje  \odkaz{trestné body}{trestný bod} a  
dále vede průběh zápasu ve spolupráci s~\odkaz{konferenciérem}{konferenciér}. 
Pro některá svá nepopulární rozhodnutí bývá terčem hodu \odkaz{papučí}{papuče} či míčků od diváků. 
 
 
Způsoby, kterými může rozhodčí ovlivňovat běh zápasu 
\begin{itemize}
\item Pískání \odkaz{faulů}{faul} a přidělování \odkaz{trestných bodů}{trestný bod}
\item Vyloučení z~následující kategorie/kategorií (např. za opakovanou chybu)
\item Výběr témat (\odkaz{Askfor}{askfor})
\item Výběr \odkaz{kategorií}{kategorie}
\item Zadávání dodatečných úkolů \odkaz{hráčům}{hráč}
\item Zastavení improvizace a např. hlasování diváků o~dalším postupu
\item Úkolování \odkaz{hudebníka}{hudebník}, \odkaz{pomocných rozhodčí}{pomocný rozhodčí} a \odkaz{konferenciéra}{konferenciér}
\item Další verbální a nonverbální výlevy
\end{itemize}
 
\subsection{ Kam dál } \begin{itemize}
\item  \odkaz{Manuál rozhodčího}{manuál rozhodčího}
\end{itemize}
 
 
 
\needspace{5cm} \section{Trenér} \label{trenér} Trenérem je myšlen jedinec, soustavně pracující s~\odkaz{týmem}{tým} na jeho zdokonalování v~improvizaci na pravidelných 
\odkaz{tréninzích}{trénink}. 
Ve starším pojetí \odkaz{zápasů}{zápas} byl rovněž v~průběhu zápasu přítomen na scéně a koučoval hráče. 
  
